\documentclass[10pt,a4paper]{article}
\usepackage[margin=1.9cm]{geometry}
\usepackage{amsmath,amssymb}
\usepackage{booktabs}
\usepackage{graphicx}
\usepackage[T1]{fontenc}
\usepackage{microtype}
\usepackage[colorlinks,linkcolor=black,citecolor=black,urlcolor=blue]{hyperref}
\usepackage{titlesec}
\titlespacing*{\section}{0pt}{6pt}{2pt}
\newcommand{\code}[1]{\texttt{#1}}
\newcommand{\op}[1]{\texttt{#1}}
\title{\vspace{-1.4cm}\textbf{What a fibre is in this matrix, and what it is not}}
\author{}
\date{\vspace{-0.9cm}\small \texttt{prototype/ecm/block\_bounce.py} $\cdot$
\texttt{prototype/ecm/block\_metrics.py} $\cdot$
\texttt{log/okuda\_ECM/02b--02h} $\cdot$ 10 August 2026}
\setlength{\parindent}{0pt}\setlength{\parskip}{0.55em}
\begin{document}\maketitle\vspace{-1.0cm}

\section*{1.\; The question, which is about resolution and not about collagen}

\op{seed\_ecm} lays the matrix down as strands: $1500$ straight segments, each carrying $60$ of the
$90{,}000$ particles, with a shared orientation and a little across-strand jitter. The movies then show
fibres being dragged and splayed, and every statement anyone has been tempted to make from those movies
is about \emph{fibres}. The operator's own docstring warns that they are geometry and not mechanics.
This note measures whether that warning is right, and it is: the strands are an arrangement of the mass,
and the material responds as though they were not there.

Two facts decide it before anything moves, and both are arithmetic on the spec. The background grid is
$64^3$, so a cell is $1/64 = 0.0156$ box units. A strand is $0.09$ long --- $5.8$ cells --- and its
particles sit $0.0015$ apart, a \emph{tenth} of a cell, while the jitter that gives the strand its
thickness is $0.004$, a \emph{quarter} of a cell. So the grid resolves a strand along its length and
cannot resolve it across: about ten particles of one strand land in a single cell and are interpolated
into one velocity. And no force is ever computed between two particles of the same strand --- MLS-MPM
computes nothing between two points directly \cite{sulsky1994,hu2018}, and \op{seed\_ecm} writes each
particle's unit fibre vector to an attribute on the operator object that no other operator reads.

A fibre in this model is therefore three things it should not be at once: a line the renderer draws by
index arithmetic (\code{particle i} belongs to \code{fibre i//per}), a clump of mass the grid smears,
and a direction stored where nothing can use it.

\section*{2.\; A load that arrives}

Every earlier run loads the matrix the same way --- slowly, outward, everywhere at once --- so none of
them can say whether it carries stress the way a fibrous material should, because there is no moment at
which the load appears. A block falling under gravity onto the floor is the cheapest load that
\emph{arrives}: the impact is a front travelling up through the strands, and the block either returns
the drop or keeps it.

Seven runs, all $1500$ strands in a cube of $0.4\times0.32\times0.4$ box units with a spherical cavity
of radius $0.07$, $720$ frames at $\Delta t = 3.2\times10^{-3}$, differing in one number apiece. What is
measured is in \code{block\_metrics.py}: the bulk (height, energy, the fraction of the drop each impact
returns), the shape (the three dimensions against their own seeded values), and the fibre state --- the
three things a strand can do that a continuum cannot, which are stretch along itself, bow sideways, and
turn.

\textbf{End-to-end, not contour.} The jitter is $0.004$ across a strand whose particles are $0.0015$
apart, so the contour is a random walk at the particle scale and is six times the end-to-end distance
before anything has moved ($0.163$ at frame $0$, and $0.163$ again at frame $360$). Any per-segment
measure reports the seeding noise. End-to-end length, and the RMS distance of a strand's particles from
its own end-to-end axis, are both taken over the whole strand, where the jitter averages down.

\section*{3.\; The strands do nothing a continuum would not}

Through an impact that compresses the block to $0.914$ of its seeded height, a strand's end-to-end
length falls to $0.982$ of the $0.0896$ box units it was seeded with, and returns to $0.999$. Its bow
goes $0.00824 \to 0.00827 \to 0.00826$ box units --- an excursion of $0.4\%$ --- and the orientation
order parameter about the drop axis moves by $0.04$. Nothing buckles, nothing stays stretched, nothing
turns.

That is not a surprise and it is not a bug: it is the prediction of the constitutive law. \op{mpm\_scatter}
computes a fixed-corotated stress from $\mathbf F$ alone and reads the strand direction nowhere, so the
only anisotropy this matrix can express is \emph{where the mass happens to be}. It is also what the
earlier alignment campaign found the expensive way --- \code{align 0.85} about one axis bought a
directional pressure difference of $1.50\times$, and a shape only appeared once the fibres were
\emph{redistributed} to put three times as many within $55^\circ$ of the poles, which changes the density
the grid sees rather than the direction it responds to.

\section*{4.\; The readout had been the wrong quantity}

\op{ecm\_stress} defaulted to \code{measure: vol}, which is $|J-1|$, the local volume change --- and an
impact is a shear. With \code{store\_stress: true} on \op{mpm\_scatter} the solver keeps the Cauchy
stress it already computes to build the grid forces, and \code{measure: vonmises} reads its invariant
instead. On the \emph{same trajectory} --- the positions are bit-identical, checked element by element
over $721\times90{,}000\times3$ floats, so this is a readout and not a mechanism --- the two differ by a
factor of $250$: at frame $200$, $|J-1|$ gives mean $0.0040$ and p99 $0.0173$, the von Mises gives mean
$1.11$ and p99 $4.16$ in stress units, peaking at p99 $42.4$ just after first contact. It costs $45\%$
more wall clock and it is the difference between seeing the front and not. One thing to carry: the
run-time \code{stress\_scale} of $2\times10^{-5}$ was chosen for $|J-1|$ and saturates every band under
von Mises; the movie re-bands from the recorded raw values (full scale $9.406$, the p99 over the run),
but \code{node\_type} in the trajectory is pinned.

\section*{5.\; The block rings, and an extracellular matrix does not}

After the impact every shape metric oscillates. The measured period is $24.2$ frames, $0.0775$ s,
$12.9$ Hz, against the $17.1$ Hz that a free--free bar of length $0.32$ at $c=\sqrt{E/\rho}=10.95$ box
units per second would have --- so it is the block's own elastic breathing mode, real physics of the
material as specified. It is nonetheless biologically wrong. The stock \code{drag} of $0.05$ is a
velocity decay time of $1/\code{drag} = 20$ s, which is $260$ ringing periods and $8.7$ times the whole
$2.31$ s run: this matrix is effectively undamped. A hydrated matrix dissipates by solvent forced
through the network, and at $\mathrm{Re}\sim10^{-10}$ it cannot ring at all.

\code{drag} is the right lever and needs no new operator --- it is a Stokes term on the network
velocity, which is what Darcy drag is. Raising it costs nothing:

\smallskip
{\small\begin{tabular}{@{}lccccc@{}}\toprule
run & \code{drag} & ripple, as a fraction & height returned at & bounding volume & median residual\\
    &             & of the seeded height  & the first impact   & at the last frame & displacement (box units)\\\midrule
\op{02c} & $0.05$ & $0.00813$ & $0.707$ & $1.172$ & $0.03604$\\
\op{02d\_drag2} & $2$ & $0.00475$ & $0.293$ & $0.999$ & $0.00264$\\
\op{02d\_drag8} & $8$ & $0.00054$ & $0.048$ & $0.997$ & $0.00108$\\
\bottomrule\end{tabular}}
\smallskip

The damping that removes the ringing removes the bounce with it: at \code{drag}$=8$ the block returns
$4.8\%$ of what it fell. There is no bouncing experiment in the regime where the material is physical,
and the three impacts of \op{02b} are a property of an undamped solid rather than of a matrix.

\section*{6.\; What was read as creep, and was not}

The undamped runs showed the block slowly inflating --- bounding volume rising to $1.34$ of its seeded
value by frame $560$, median particle displacement with the block's own translation removed reaching
$0.036$ box units, which is $2.3$ grid cells --- and I reported that as the one irreversible thing the
material did. It is not the material. At \code{drag}$=2$ the volume ends at $0.999$ and the residual
displacement at $0.0026$, smaller by a factor of $13.6$; at \code{drag}$=8$, $0.0011$. It was undamped
particle noise accumulating, and the only reason it looked like creep is that it was monotone.

The general form of that mistake is worth keeping: a quantity that only ever increases is not thereby a
material property, and the cheapest test of whether it is one is to change the damping, which no
material property should notice.

\begin{figure}[t]\centering
\includegraphics[width=\textwidth]{../log/okuda_ECM/02c_ecm_block_vonmises/material.png}
\caption{\op{02c}, the undamped block, $721$ frames. \textbf{a--c} three impacts, and the height returned
falling $0.707\to0.612\to0.416$, which reads as dissipation and is in fact the block coming apart ---
\textbf{e} shows its bounding volume growing to $1.34$ of the seeded value. \textbf{d--e} ring at
$12.9$ Hz, the block's own elastic mode. \textbf{f--g} are the fibre state and are the point: the
end-to-end stretch recovers to $0.999$ and the bow moves by $0.4\%$, so the strands ride the continuum.
Everything in \textbf{e}, and the orange trace in \textbf{h}, is removed by \code{drag}$=2$.}
\end{figure}

\section*{7.\; What each lever costs, measured}

Four runs, each changing one number from the \code{drag}$=2$ configuration, all $720$ frames, solve time
timed inside the script and excluding rendering. The last row changes three numbers at once and is the
configuration \S9 proposes.

\smallskip
{\small\begin{tabular}{@{}llccrrcccc@{}}\toprule
run & what changed & particles & substeps & solve (s) & speed-up &
first impact & height returned & peak compression & fibre stretch\\
 & & & per frame & & & (frame) & at that impact & (of seeded height) & (minimum)\\\midrule
\multicolumn{10}{@{}l}{\textit{against the \code{drag}=2 control}}\\
\op{02d\_drag2} & --- & $90{,}000$ & $16$ & $\sim678^\ast$ & $1.0$ & $181$ & $0.293$ & $0.9468$ & $0.9881$\\
\op{02e\_sub4} & $\Delta t_{\rm sub}\,4\!\times\!10^{-4}$ & $90{,}000$ & $8$ & $339.1$ & $2.0$ & $181$ & $0.304$ & $0.9439$ & $0.9880$\\
\op{02f\_sub8} & $\Delta t_{\rm sub}\,8\!\times\!10^{-4}$ & $90{,}000$ & $4$ & $180.6$ & $3.8$ & $182$ & $0.314$ & $0.9415$ & $0.9879$\\
\op{02g\_p30k} & $20$ particles per strand & $30{,}000$ & $16$ & $261.9$ & $2.6$ & $181$ & $0.280$ & $0.9443$ & $0.9875$\\\midrule
\multicolumn{10}{@{}l}{\textit{against the \code{drag}=8 control}}\\
\op{02d\_drag8} & --- & $90{,}000$ & $16$ & $\sim678^\ast$ & $1.0$ & $314$ & $0.048$ & $0.9868$ & $0.9962$\\
\op{02h\_lean} & all three & $30{,}000$ & $8$ & $\mathbf{115.5}$ & $\mathbf{5.9}$ & $315$ & $0.043$ & $0.9854$ & $0.9960$\\
\bottomrule\end{tabular}}

{\footnotesize $^\ast$ implied, not timed: the timer was added after those runs. Solve time is exactly
linear in substeps at fixed particle count ($339.1$ s at $8$, $180.6$ s at $4$, i.e.\ $42.4$ and $45.2$ s
per substep-set), and \op{02c}'s wall clock including compression was $1086$ s.}
\smallskip

Two levers, and both were being paid for without buying anything. The substep was $\Delta t_{\rm sub} =
2\times10^{-4}$ against a CFL limit of $\mathrm{d}x/c = 1.4\times10^{-3}$ --- a factor of seven of
headroom. Halving the substep count moves the peak compression by $0.3\%$ and the fibre stretch by
$0.01\%$; quartering it moves them by $0.6\%$ and $0.02\%$ and is still stable, which is why the
proposal takes the safe half rather than the quarter. And $60$ particles per strand put ten of them in
one grid cell, so two thirds of them are describing a strand at a resolution the grid discards: at $20$
per strand the peak compression moves by $0.26\%$ and the run is $2.6\times$ cheaper. Combined, the
lean configuration reproduces its control to within one frame on the impact, $0.14\%$ on the peak
compression and $0.02\%$ on the fibre stretch, at $5.9\times$ less compute.

\section*{8.\; What would falsify a fibre model}

Nothing above needs fibres to explain it, so nothing above can falsify a fibre model either. What a real
fibre network does, in the order the property matters: it carries load along its axis in tension and
buckles under $\sim\pi^2EI/L^2$ in compression, which for a fibril of aspect ratio $100$ is $10^{-4}$ of
its tensile stiffness; that asymmetry is what produces strain stiffening as the network goes from
bending-dominated to stretch-dominated near $10$--$20\%$ strain \cite{storm2005,licup2015}, and
long-range transmission, where displacement decays like $r^{-0.5}$ to $r^{-1}$ instead of the $r^{-2}$
of a linear continuum \cite{wang2014}; and its fibres rotate toward the maximum principal stretch, so
they end up circumferential around a compressed region.

Three loadings would show it, in increasing cost:

\begin{itemize}\setlength{\itemsep}{1pt}
\item \textbf{Compress between plates to $20\%$ and release.} The tangent modulus at $5\%$ against
$20\%$ strain is the stiffening ratio, which a network raises by one to two orders of magnitude and an
isotropic solid does not raise at all; the residual strain after unloading is the plasticity; and the
effective Poisson ratio separates a solid that bulges ($\approx0.3$ from the Lam\'e pair) from a network
that densifies ($0$ to $0.1$).
\item \textbf{Pull by $10\%$, push by $10\%$.} The ratio of secant moduli is $1$ by construction for the
present law and must exceed $3$ for anything that buckles \cite{notbohm2015}.
\item \textbf{Contract a sphere at the centre by $10\%$} and fit $u(r)\propto r^{-n}$. An incompressible
linear continuum gives $n=2$, a fibrous matrix $n$ between $0.5$ and $1$ \cite{wang2014}. That single
exponent is the property the whole spheroid--matrix coupling rests on, and it is the test to run first.
\end{itemize}

One metric belongs in every one of them: the mean $\cos^2\theta$ between a strand's axis and the
maximum-principal-stress eigenvector of its own particles' $\boldsymbol\sigma$. Random gives $1/3$; a
fibre-reinforced law drives it toward $1$ wherever the strand is in tension. That is ``builds stress
along its direction'' written as a number that can fail, and it needs the one structural change
everything else waits on: the fibre direction as a per-particle buffer rather than an attribute on the
operator.

\section*{9.\; The spec I propose for the ECM--BM--spheroid study}

Keep the two passes. The reason is unchanged and is geometric, not a preference: the MPM grid is fixed
to the unit cube, the epithelium lives in a box fifty times larger, and shrinking the tissue to fit
changes which term of its energy dominates. What crosses between them is the \op{surface} Level.

\textbf{Pass 1, the epithelium}, carries no MPM and costs minutes: $200\to\sim\!1100$ cells over $402$
frames, \op{shape\_energy\_3d} / \op{reconnect\_t1\_3d} / \op{divide\_3d}, with
\op{ecm\_growth\_gate\_3d} reading the recorded pressure map \cite{helmlinger1997,montel2011}. Nothing
in this note touches it.

\textbf{Pass 2, the matrix and the sheet}, is where the compute is, and where the four numbers above go:

\smallskip
{\footnotesize
\begin{verbatim}
  fields:  mpm_grid: {n_grid: 64}            # dx = 0.0156; a strand is 5.8 cells long
  sets:
    stroma:       {parent: ecm_body,  per_parent: 40000, types: {ecm: {youngs: 15}}}
    bm:           node set, NOT mpm_particle    # see below
    plaque:       edge set  cell -> bm
    fibril:       edge set  bm   -> ecm
  operators:
    - {op: seed_ecm,     at: stroma, n_fibres: 6000, fibre_len: 0.16, per_strand: 20}
    - {op: mpm_scatter,  at: stroma, drag: 8.0, store_stress: true}   # overdamped, keeps sigma
    - {op: ecm_stress,   at: stroma, measure: vonmises, scale: <p99 of a calibration frame>}
  schedule:
    - {substep_dt: 4.0e-4, steps: [mpm_strain, mpm_scatter, mpm_grid_update, mpm_gather]}
\end{verbatim}}
\smallskip

Four choices, each with the measurement behind it. \code{drag: 8} because at $0.05$ the matrix rings for
the entire run and accumulates $2.3$ grid cells of displacement that is not material (\S5, \S6), and
raising it costs nothing. $\Delta t_{\rm sub} = 4\times10^{-4}$, eight substeps rather than sixteen,
because the CFL limit is $1.4\times10^{-3}$ and halving the count moves every deformation measure by
less than $0.3\%$ (\S7). Twenty particles per strand rather than sixty, because ten of the sixty land in
one grid cell (\S1) and dropping to twenty moves the same measures by less than $0.3\%$ --- for the
spheroid's $6000$ strands that is $40{,}000$ particles, not $120{,}000$. And
\code{store\_stress}/\code{vonmises} because $|J-1|$ cannot see a shear (\S4). Together these are the
$5.9\times$ of \op{02h}: a pass-2 of about ninety seconds rather than ten minutes, which is what makes a
sweep affordable rather than a weekend.

\textbf{The basement membrane must not be an MPM material, and the reason is a number.} A basement
membrane is $\sim\!100$ nm thick on a spheroid $\sim\!200\,\mu$m across --- one part in two thousand of
the radius, which at a spheroid radius of $0.2$ box units is $10^{-4}$ box, or one part in $150$ of a
grid cell. There is no resolution at which an MPM sheet is a sheet here. Worse, it would set the
timestep for everything else: the sheet is $10$ to $100$ times stiffer than the stroma, and the substep
scales as $1/\sqrt{E}$, so \code{youngs} $=100\times$ the stroma's takes the run from three substeps per
frame to twenty-seven --- \emph{nine times the cost, to resolve a layer the grid cannot represent}. The
sheet therefore stays what \S6 of \code{spheroid\_ecm} argues it is: nodes and crosslinks, massless and
overdamped, integrated at its own rate, with \op{plaque} as an edge-set so the adhesion returns its
reaction to the cell \cite{kanchanawong2010} and \op{fibril} as an edge-set carrying the stroma's load
through the sheet \cite{keene1987} rather than letting the epithelium push the stroma directly.

\textbf{Then spend the saving on the one thing that changes an answer.} In order, and each is a separate
operator with its own falsifier: (i) the fibre direction as a per-particle buffer, which nothing else
can proceed without; (ii) a tension-only transversely isotropic term in the stress, $I_4 = |\mathbf F
\mathbf a|^2$ active only when $I_4>1$ \cite{hgo2000,gasser2006}, which is the only way stress can align
with a strand, tested by the $\cos^2\theta$ metric of \S8 and by the $r^{-n}$ exponent; (iii) affine
reorientation $\mathbf a \leftarrow \mathbf F\mathbf a/|\mathbf F\mathbf a|$ \cite{wang2014}; (iv) mass
balance on the sheet rather than a fixed particle budget \cite{matsubayashi2020}. Items (ii) and (iii)
are perhaps a $15\%$ surcharge on pass 2 --- an $I_4$ and a rank-one stress per particle per substep ---
against the $5.9\times$ already banked.

\textbf{One caveat that the saving makes sharper.} Under one-way coupling the stroma's \code{youngs} is
very nearly a free parameter: the tissue never feels the matrix, and the growth gate normalises the
pressure map to its own late-run median, so a softer stroma changes the substep it needs and not the
shape that comes out. That is convenient --- softer is cheaper, since the substep scales as
$1/\sqrt{E}$ --- and it is also the clearest statement of what one-way coupling costs. The moment an
operator writes back to \op{surface}, \code{youngs} stops being free, and the numbers in this section
have to be re-measured rather than carried over.

\section*{10.\; References}

{\small
\begin{thebibliography}{99}\setlength{\itemsep}{0pt}
\bibitem{sulsky1994} Sulsky, D., Chen, Z., Schreyer, H.L. (1994) A particle method for history-dependent
materials. \emph{Comput. Methods Appl. Mech. Eng.} \textbf{118}:179--196.
\bibitem{hu2018} Hu, Y. \emph{et al.} (2018) A moving least squares material point method with
displacement discontinuity and two-way rigid body coupling. \emph{ACM Trans. Graph.} \textbf{37}(4):150.
\bibitem{storm2005} Storm, C., Pastore, J.J., MacKintosh, F.C., Lubensky, T.C., Janmey, P.A. (2005)
Nonlinear elasticity in biological gels. \emph{Nature} \textbf{435}:191--194.
\bibitem{licup2015} Licup, A.J. \emph{et al.} (2015) Stress controls the mechanics of collagen networks.
\emph{Proc. Natl. Acad. Sci. USA} \textbf{112}(31):9573--9578.
\bibitem{notbohm2015} Notbohm, J., Lesman, A., Rosakis, P., Tirrell, D.A., Ravichandran, G. (2015)
Microbuckling of fibrin provides a mechanism for cell mechanosensing. \emph{J. R. Soc. Interface}
\textbf{12}(108):20150320.
\bibitem{wang2014} Wang, H., Abhilash, A.S., Chen, C.S., Wells, R.G., Shenoy, V.B. (2014) Long-range
force transmission in fibrous matrices enabled by tension-driven alignment of fibers.
\emph{Biophys. J.} \textbf{107}(11):2592--2603.
\bibitem{hgo2000} Holzapfel, G.A., Gasser, T.C., Ogden, R.W. (2000) A new constitutive framework for
arterial wall mechanics and a comparative study of material models. \emph{J. Elasticity} \textbf{61}:1--48.
\bibitem{gasser2006} Gasser, T.C., Ogden, R.W., Holzapfel, G.A. (2006) Hyperelastic modelling of arterial
layers with distributed collagen fibre orientations. \emph{J. R. Soc. Interface} \textbf{3}(6):15--35.
\bibitem{helmlinger1997} Helmlinger, G., Netti, P.A., Lichtenbeld, H.C., Melder, R.J., Jain, R.K. (1997)
Solid stress inhibits the growth of multicellular tumor spheroids. \emph{Nat. Biotechnol.}
\textbf{15}(8):778--783.
\bibitem{montel2011} Montel, F. \emph{et al.} (2011) Stress clamp experiments on multicellular tumor
spheroids. \emph{Phys. Rev. Lett.} \textbf{107}:188102.
\bibitem{kanchanawong2010} Kanchanawong, P. \emph{et al.} (2010) Nanoscale architecture of
integrin-based cell adhesions. \emph{Nature} \textbf{468}:580--584.
\bibitem{keene1987} Keene, D.R., Sakai, L.Y., Lunstrum, G.P., Morris, N.P., Burgeson, R.E. (1987)
Type VII collagen forms an extended network of anchoring fibrils. \emph{J. Cell Biol.}
\textbf{104}(3):611--621.
\bibitem{matsubayashi2020} Matsubayashi, Y. \emph{et al.} (2020) Rapid homeostatic turnover of embryonic
ECM during tissue morphogenesis. \emph{Dev. Cell} \textbf{54}(1):33--42.
\end{thebibliography}}

\end{document}
